% Experiment report. Fixed sections: see experiments/TEMPLATE/report.tex and experiments/README.md.
% Build: npm run experiment:report -- 21-rl-strategic-vi-perf
\documentclass{../polyexp}
\begin{document}
\input{generated}
\ExpHeader

\section{Question and hypothesis}
\textbf{The idea (David, 26 Nov 2026, issue \#166).} Repeat experiment 15 on the Perfection-tuned bot of experiment
17. In 15 the base bot was tuned for Domination, and most of the policy's gain came from ``research more'', which
looks like fixing a mode mismatch rather than strategy. With \texttt{greedy-full-probe-perf} as the base bot, the
question is whether value iteration over the six integers (turn, cities, income, units, sum of city levels, techs)
still finds something the tuned bot misses, for example playing differently early and late.

\textbf{Prediction} (written before any run of this experiment, including the action check):
\begin{enumerate}
  \item \emph{Less to gain from a fixed action.} In the action check (2{,}000 seeds) the best action played every
        turn gains less than half of 15's research gain against the bot: under 225 points in the mean
        (15: +450).
  \item \emph{The policy still beats the bot, by less.} On the 10{,}000 evaluation seeds the policy passes the
        \#171 rule against \texttt{greedy-full-probe-perf}, with a mean gain of 30--200 points (15: +118 against the
        untuned bot).
  \item \emph{Values are optimistic.} $V$(start) at min-count 1 exceeds the policy's real mean final score minus
        the start score by more than 500 points; the gap shrinks with a higher min-count (15: +4{,}202 at $m=1$).
  \item \emph{No better than the best fixed action.} Against the best action of the check played every turn, the
        policy does not pass the \#171 rule (as in 15).
\end{enumerate}
\textbf{Against the idea:} if the policy fails against the tuned bot, strategic boosts over this state add nothing
once the bot is tuned for the mode; experiment 15's gain was the mode mismatch.

\section{Mechanism}
\textbf{Strategic actions.} At the start of each learner turn one action $a \in \{0,\dots,4\}$ is chosen and holds
for the whole turn: it multiplies the bot's move-type weights (\texttt{FullWeights::ty}, the factor on the score of
every move of that type) by the boost 2. 0 \emph{bot}: no change; 1 \emph{expand}: capture, examine;
2 \emph{economy}: build; 3 \emph{army}: train, attack; 4 \emph{research}: research
(\texttt{engine/rlvalue/src/vi.rs}, \texttt{BOOSTED}).

\textbf{Data.} 300{,}000 games in which the learner picks a uniformly random action each turn (fixed by seed and
turn); each row records the state, the score and the action.

\textbf{Value iteration.} For each state $s$ and action $a$ the observed transitions $(r, s')$ form $\hat P(s'\mid
s,a)$ and $\hat r$. Backwards from the last turn: $Q(s,a) = \frac{1}{n(s,a)}\sum \bigl(r + V(s')\bigr)$,
$V(s) = \max_{a:\,n(s,a) \ge m} Q(s,a)$, $V(\text{terminal}) = 0$; the policy is the argmax. A state where no action
has $m$ visits keeps the mean over all its transitions and has no policy entry. When playing, a state without an
entry gets action 0 (the bot).

\textbf{Choice of $m$.} $m \in \{1, 5, 20, 50\}$; each policy plays 2{,}000 validation games (seeds
$700{,}000{,}000\dots$) against the bot on the same seeds; the $m$ with the highest share of better seeds is kept.
Only that one is evaluated on the evaluation seeds.

\section{Setup}
As experiment 15 except the setup, which is now the decided gate setup on \texttt{dev} at da04376: Perfection,
16$\times$16 Continents, 4 players, Crazy, seat 0 Xin-Xi as the human seat; the learner's base bot
\texttt{greedy-full-probe-perf}; opponents \texttt{greedy-full-probe-perf}, \texttt{greedy-full-combo-perf} and the
anchor \texttt{greedy-full-probe-perf}, shuffled over the seats per seed by the gate's \texttt{seat\_assignment}.
Seeds: training $1\dots300{,}000$; action check $600{,}000{,}000\dots$ (2{,}000); validation
$700{,}000{,}000\dots$ (2{,}000); evaluation $10^9\dots$ (10{,}000).

\textbf{Verdict tool.} \texttt{rl-gate} cannot play a policy yet; the learner seam is being added (DECISION of 26 Nov
on \#166, own PR). Until it is merged, \texttt{rl-value eval} plays the policy and the base on the same seeds and
applies the \#171 rule itself (labelled as such); \texttt{rl-value}'s games are identical to \texttt{rl-gate}'s
(500 of 500 seeds, experiment 20). By the pass rule of 26 Nov an RL version is judged against the previous RL
version, so the comparison with \texttt{greedy-full-probe-perf} here is the baseline comparison, and this is the
candidate for RL v1 of David's track.

\textbf{Action check} (each action every turn, 2{,}000 seeds $600{,}000{,}000\dots$, \texttt{check-boost2.json},
\texttt{check-boost4.json}; run after the predictions above were written):

\begin{center}\small
\begin{tabular}{lrrrrr}
\toprule
action (boost 2) & mean score & median & 10th pct & techs at end & better / worse vs bot \\
\midrule
bot      & 6871 & 8080 & 1420 & 18.49 & -- \\
expand   & 6887 & 8128 & 1420 & 18.50 & 84 / 61 \\
economy  & 6654 & 7852 & 1380 & 18.27 & 760 / 886 \\
army     & 6861 & 8068 & 1400 & 18.55 & 404 / 426 \\
research & 7276 & 8845 & 1464 & 19.78 & 1{,}098 / 731 \\
\bottomrule
\end{tabular}
\end{center}
With boost 4, research gains less (6{,}929) and costs cities (2.19 against 2.61) and income, economy loses more
(6{,}377), expand and army change little; boost 2 is kept, as in 15. The best fixed action is research
(\texttt{fixed:4}), so the final comparison against a fixed action is against ``research every turn''.

Commands (from \texttt{engine/}):
\begin{verbatim}
B=./target/release/rl-value
$B check --games 2000 --seed-start 600000000 --boost 2 --out ../experiments/21-rl-strategic-vi-perf/check-boost2.json
$B check --games 2000 --seed-start 600000000 --boost 4 --out ../experiments/21-rl-strategic-vi-perf/check-boost4.json
$B collect --driver random --games 300000 --seed-start 1 --out target/rl-value/p4-train-random.tsv
$B fit-vi --train target/rl-value/p4-train-random.tsv --min-count 1,5,20,50 --out target/rl-value/vi-perf
$B eval --policy target/rl-value/vi-perf/policy-mM.tsv --n 2000 --seed-start 700000000 --out .../val-mM.json
$B eval --policy <chosen> --n 10000 --out .../eval-vs-bot.json
$B eval --policy <chosen> --n 10000 --base-driver fixed:K --out .../eval-vs-research.json   # K = 4, research
../engine/target/release/rl-gate --candidate rl-policy:experiments/21-rl-strategic-vi-perf/policy-m5.tsv \
  --base greedy-full-probe-perf --out ../experiments/21-rl-strategic-vi-perf
\end{verbatim}

\section{Results}
\ExpResults

All numbers are copied from \texttt{check-boost*.json}, \texttt{vi-summary.json}, \texttt{val-m*.json},
\texttt{eval-vs-bot.json}, \texttt{eval-vs-research.json} and the \texttt{rl-gate} JSON in this folder (26 Nov
2026). Data: 300{,}000 random-action games, 8{,}493{,}078 rows, 506{,}245 distinct states; collection 1{,}544\,s on 10
threads. The chosen policy is \texttt{policy-m5.tsv} (its header names the base bot and the boost).

\paragraph{Value iteration and validation} (2{,}000 validation seeds, policy against the bot on the same seeds;
the bot scores 6{,}840 there; optimism = $V$(start) minus the policy's real mean remaining score, start score 515):

\begin{center}\small
\begin{tabular}{rrrrrrr}
\toprule
min-count & policy covers & $V$(start) & optimism & policy mean & better / worse & better share \\
\midrule
1  & 100.0\% & 10{,}881 & $+4{,}427$ & 6{,}969 & 867 / 733 & 54.2\% \\
5  & 86.9\%  & 8{,}272  & $+1{,}793$ & 6{,}994 & 853 / 701 & \textbf{54.9\%} \\
20 & 71.9\%  & 7{,}421  & $+968$     & 6{,}968 & 785 / 668 & 54.0\% \\
50 & 58.5\%  & 7{,}016  & $+581$     & 6{,}950 & 715 / 614 & 53.8\% \\
\bottomrule
\end{tabular}
\end{center}
Min-count 5 is kept (the highest share of better seeds). On validation all four fail the \#171 rule only on the
10th-percentile guard (1{,}410--1{,}420 against 1{,}440); the sign test rejects for all four ($p \le 0.003$).

\paragraph{Final test} (10{,}000 evaluation seeds $10^9\dots$; \texttt{rl-gate} with
\texttt{--candidate rl-policy:\dots/policy-m5.tsv}, from PR \#198, not yet on \texttt{dev}; \texttt{rl-value eval}
gives the same numbers seed for seed):

\begin{center}\small
\begin{tabular}{lrrrrrrl}
\toprule
against & policy mean / median / p10 & other mean / median / p10 & better / worse / ties & mean $d$ (99\% bounds) & rule \\
\midrule
bot (probe-perf) & 6943 / 8350 / 1415 & 6839 / 8112 / 1419.5 & 4277 / 3536 / 2187 & $+104$ ($+57$, $+152$) & FAIL (p10) \\
research every turn & 6943 / 8350 / 1415 & 7295 / 8885 / 1485 & 3380 / 5430 / 1190 & $-352$ (upper $-312$) & FAIL \\
\bottomrule
\end{tabular}
\end{center}
Against the bot the sign test rejects clearly (54.74\% of 7{,}813 differing seeds better, 99\% lower bound 53.42\%,
$p = 2.7\cdot10^{-17}$) and the mean guard holds, but the 10th percentile is 4.5 points lower (1{,}415 against
1{,}419.5), so the rule says FAIL. Win rate 3.3\% against 3.1\%. The information-only run with 15 opponents (5{,}000
seeds) shows no gain: mean 2{,}579 against 2{,}679, better on 1{,}408 seeds and worse on 1{,}430.

\section{Behaviour}
On the evaluation seeds 86.7\% of learner turns are in a state with a policy entry. Actions over all learner turns:
bot 28.1\%, expand 15.1\%, economy 12.2\%, army 15.2\%, research 29.5\%. The policy picks research most, but only on
three turns in ten, and picks economy on one turn in eight although economy every turn costs 217 points in the
check; expand and army, which change almost nothing when played every turn, take 30\% of turns together.

\section{Interpretation}
\textbf{Prediction 1 (best fixed action gains under 225 points): wrong.} Research every turn gains 405 points in
the check and 456 on the evaluation seeds (7{,}295 against 6{,}839), about as much as against the untuned bot in 15
(+450). The Perfection fine-tune of experiment 17 lowered the bot's research factor further (to 0.417); doubling it
still pays. So the tuned bot still researches too little for this setup. One untested guess why: the trainer's
games have no human seat and two sides on alternating seats.

\textbf{Prediction 2 (passes the \#171 rule against the bot, gain 30--200): half right.} The gain is in the
predicted range (+104 in the mean, 54.7\% of differing seeds better, very significant), but the policy fails the
rule on the 10th-percentile guard by 4.5 points: the policy does not protect the bad games (eliminations), where
its actions make no difference or slightly hurt.

\textbf{Prediction 3 (values optimistic, the gap shrinks with min-count): right} (+4{,}427 at $m=1$ down to +581 at
$m=50$).

\textbf{Prediction 4 (not better than the best fixed action): right}, and by a lot: research every turn beats the
policy on 5{,}430 seeds against 3{,}380, by 352 points in the mean.

\textbf{Why value iteration does not find ``research every turn''.} A policy that always researches is in the
policy class (one action per state), and it is clearly better, yet value iteration picks research on only 30\% of
turns. The Behaviour section points to noise, not to a real trade-off: the policy spreads its choices over expand,
army and economy, which played every turn are no better or worse than the bot. With five actions and about 17
visits per state, each $Q(s,a)$ rests on a handful of transitions; taking the maximum over five noisy estimates
picks whichever action was lucky in that state (the same bias that makes $V$(start) optimistic), and the backups
spread that bias to earlier turns. The state is also not Markov (experiment 20), so successor values pool games
that differ in what the six integers do not show. Experiment 15's conclusion carries over to the tuned bot: on this
state the learned policy beats the bot it steers on most seeds, but a fixed rule found by a five-action check beats
it.

\todo{david}{Your reading: agree, disagree, what is missing?}

\section{Threats to validity}
\begin{itemize}
  \item \emph{Random-action data.} The model is learned from uniformly random actions; the states a policy that
        researches more visits are rarer in the data than in its own play (86.7\% of its turns have an entry).
  \item \emph{The p10 guard decides by 4.5 points.} The 10th percentile sits in the elimination games; a slightly
        different seed set could flip the verdict either way. The sign test and the mean are far from the edge.
  \item \emph{Gate code not on \texttt{dev} yet.} The \texttt{rl-gate} run uses the learner seam of PR \#198; it agrees
        seed for seed with \texttt{rl-value eval}, and an empty policy ties its base bot on all seeds tested.
  \item \emph{One boost, one action set.} Boost 2 and these five actions; other action designs (spending decisions,
        asked on \#166) are not tested.
\end{itemize}

\section{Verdict and next step}
\textbf{Against the greedy baseline: FAIL} (information under the pass rule of 26 Nov: greedy models do not decide
an RL pass). The policy is better than \texttt{greedy-full-probe-perf} on significantly more seeds (+104 in the
mean) but fails the 10th-percentile guard by 4.5 points, and research every turn beats it by 352 points.
Whether \texttt{policy-m5} becomes RL v1 of David's track (v1 enters by decision once it plays complete games) is
David's call.

Next step ideas (one per experiment): (a) fewer actions (bot and research only), so each $Q$ rests on more
transitions; (b) double estimation against the maximisation bias (choose the action on one half of the games,
value it on the other); (c) research as the default action in states without a policy entry; (d) collect the data
with the policy itself instead of random actions.

\section{Notes for the paper}
Claim: on the Perfection-tuned bot, value iteration over the six-integer state gives a policy that is better than
the bot on 54.7\% of the differing seeds (+104 points), but a fixed rule (research every turn, +456) is far better,
because the maximum over noisy $Q$ estimates picks lucky actions. A clean example of the maximisation bias, with the
optimism of $V$(start) as the measured symptom (+4{,}427 at min-count 1, +581 at 50).

\todo{david}{Is this worth a place in the talk or report? Where?}

\end{document}
